\documentclass[11pt]{article}

\usepackage{listings}
\usepackage{xcolor}

\definecolor{codegreen}{rgb}{0,0.6,0}
\definecolor{codegray}{rgb}{0.5,0.5,0.5}
\definecolor{codepurple}{rgb}{0.58,0,0.82}
\definecolor{backcolor}{rgb}{0.95,0.95,0.95}

\lstset{
  backgroundcolor=\color{backcolor},
  commentstyle=\color{codegreen},
  keywordstyle=\color{magenta},
  numberstyle=\tiny\color{codegray},
  stringstyle=\color{codepurple},
  basicstyle=\ttfamily\footnotesize,
  breaklines=true,
  captionpos=b,
  numbers=left,
  numbersep=5pt,
  showstringspaces=false,
  tabsize=2,
  language=Java
}

\begin{document}

\centerline{\Large\bf Java Code Reference}\par
\bigskip

A basic Java class with a \lstinline{main} method:

\begin{lstlisting}
public class Hello {
    public static void main(String[] args) {
        String message = "Hello, World!";
        System.out.println(message);

        // Loop example
        for (int i = 0; i < 5; i++) {
            System.out.println("i = " + i);
        }
    }
}
\end{lstlisting}

Inline: use \lstinline{System.out.println()} to print to stdout.

\end{document}
